\subsection{Reversing the chain rule and product rule}

\begin{corebox}[Only the elementary cases]
Substitution and integration by parts are introduced to reveal the ideas behind them. The examples are chosen so that the method, not algebraic complexity, is the main difficulty.
\end{corebox}

\begin{examplebox}[A direct substitution]
Consider
\[
\int 2x(x^2+1)^3\,dx.
\]
Let
\[
u=x^2+1,
\qquad du=2x\,dx.
\]
Then
\[
\int 2x(x^2+1)^3\,dx=\int u^3\,du=\frac{u^4}{4}+C
=\frac{(x^2+1)^4}{4}+C.
\]
\end{examplebox}

\begin{interpretationbox}[Substitution]
Substitution reverses the chain rule: an inner function and its derivative are recognized as one new variable.
\end{interpretationbox}

\begin{theorembox}[Integration by parts]
\[
\int u\,dv=uv-\int v\,du.
\]
\end{theorembox}

\begin{examplebox}[An elementary example by parts]
For
\[
\int xe^x\,dx,
\]
choose $u=x$ and $dv=e^x\,dx$. Then $du=dx$ and $v=e^x$, so
\[
\int xe^x\,dx=xe^x-\int e^x\,dx=xe^x-e^x+C.
\]
\end{examplebox}

\begin{interpretationbox}[Integration by parts]
The method reverses the product rule. One factor is differentiated while the other is integrated.
\end{interpretationbox}
